% !TEX root = sga4.5-en.tex
% ENGLISH EDITION -- TRANSLATED FROM FRENCH (AI-assisted, needs human review).

\chapter{$L$-functions modulo \texorpdfstring{$\ell^n$}{l n} and modulo \texorpdfstring{$p$}{p}}\label{III}

In this expos\'e, I prove a formula analogous to the ratio theorem on
the trace formula (this volume; cited ``Rapport'' in what follows) for a
coefficient ring $A$ which is torsion prime to $p$, or a field of
characteristic $p$. Essential use will be made of \cite[XVII 5.5]{sga4}.










\section{Symmetric tensors}\label{III:1}




\subsection{}\label{III:1-1}

Let $A$ be a commutative ring. For a list of properties of the functors
$\Gamma^n:(\text{$A$-modules})\to (\text{$A$-modules})$, I refer to
\cite[XVII 5.5.1 et 2]{sga4}. Let us only recall that for $M$ flat (in fact, only
the case where $M$ is projective of finite type interests us), one has
\[
  \Gamma^n M = \left(M^{\otimes n}\right)^{S_n} \qquad \text{(symmetric tensors of degree $n$).}
\]




\subsection{}\label{III:1-2}

Let $X$ be an $S$-scheme, $X^n/S$ the $n$-th fibred power of $X$ and
$\pi$ the projection from $X^n/S$ onto $(X^n/S)/S_n=\operatorname{Sym}_S^n(X)$
(assumed to exist). If $\sG$ is a sheaf of $A$-modules on $X$, the
external tensor product of $n$ copies of $\sG$; $\sG^{\boxtimes n}$, is an
$S_n$-equivariant sheaf on $X^n/S$. We shall have to make use of the
\emph{external} functor $\Gamma^n$ of \cite[XVII 5.5.7 \`a 9]{sga4}. Let us only
recall that for $\sG$ a sheaf of flat $A$-modules (in fact, only the
constructible and flat case interests us), one has
\[
  \Gamma_\text{ext}^n(\sG) = \left(\pi_*\sG^{\boxtimes n}\right)^{S_n} \text{.}
\]




\begin{lemma_}\label{III:1-3}
For $S=\spec(k)$, with $k$ algebraically closed, $X$ a finite sum of
copies of $S$, and $G$ a flat sheaf of $A$-modules on $X$, one has
\[
  \Gamma\left(\operatorname{Sym}_S^n(X),\Gamma_\text{ext}(\sG)\right) = \Gamma^n \Gamma(X,\sG)
\]
\end{lemma_}

Indeed,
\begin{align*}
  \Gamma\left(\operatorname{Sym}_S^n(X),\Gamma_\text{ext}^n(\sG)\right)
    &= \Gamma\left(\operatorname{Sym}_S^n(X),\pi_*\sG^{\boxtimes n}\right)^{S_n} \\
    &= \Gamma\left(X^n/S,\sG^{\boxtimes n}\right)^{S_n}
    = \left(\Gamma(X,\sG)^{\otimes n}\right)^{S_n}
    = \Gamma^n \Gamma(X,\sG) \text{.}
\end{align*}




\subsection{}\label{III:1-4}

We shall have to use the derived functors of the non-additive functors $\Gamma^n$
and $\Gamma_\text{ext}^n$.

The theory of such derived functors is due to Dold and Puppe
\cite{do61}. For more complete reminders than those given below, I
refer to \cite[XVII 5.5.3]{sga4}.

Let $\cA$ be an abelian category (or even an additive category in which
every idempotent endomorphism is the projection onto a direct factor). Let
$N$ denote the normalization functor: $N:($cosimplicial complexes of objects of
$\cA)\to($differential complexes of objects of $\cA$, with $K^i$ for
$i<0)$. It is an equivalence of categories. Like its inverse, it
transforms homotopic morphisms into homotopic morphisms. If $\Gamma$ is a functor
$\cA\to\cB$, and if $K$ is a complex of objects of $\cA$, in positive
degrees ($K^i=0$ for $i<0$), one puts $T K=N T N^{-1} K$.




\subsection{Example}\label{III:1-5}

If $K$ is reduced to $M$ in degree $0$, $T K$ is reduced to $T M$ in
degree $0$.




\subsection{}\label{III:1-6}

If $K$ is a bounded complex, in positive degrees, of projective
$A$-modules of finite type, the complex of $A$-modules $\Gamma^n K$ has the same
properties.

Let $\eD_\text{parf}^{\geqslant 0}(A)$ denote the subcategory of
$\eD_\text{parf}(A)$ (Rapport, \ref{II:4-3}) formed by complexes of
tor-dimension $\leqslant 0$. One sees that $\Gamma^n$ induces a functor
\[
  \eL\Gamma^n : \eD_\text{parf}^{\geqslant 0}(A) \to \eD_\text{parf}^{\geqslant 0}(A) \text{.}
\]

For the analogous, but more complicated, definition of
$\eL\Gamma_\text{ext}^n$, I refer to \cite[XVII 5.5.14]{sga4}.




\subsection{}\label{III:1-7}

Let $K$ be a bounded complex of projective $A$-modules of finite type, and $u$
an endomorphism of $K$. We denote by $\det(1-u t,K)$ the formal power series with
constant term $1$
\[
  \det(1-u t,K) = \prod_i \det(1-u t,K^i)^{(-1)^i} \text{.}
\]

We leave to the reader the verification of the following properties.


\subsubsection{}\label{III:1-7-1}

Let $F$ be a finite filtration stable under $u$, such that the $\gr_F^p(K^q)$
are projective of finite type. Then
\[
  \det(1-u t,K) = \prod_p \det(1-u t,\gr_F^p(K)) \text{.}
\]


\subsubsection{}\label{III:1-7-2}

If one shifts $K$ by $q$ steps to the left, one has
\[
  \det(1-u t,K[q]) = \det(1-u t,K)^{(-1)^q} \text{.}
\]




\begin{proposition_}\label{III:1-8}
Let $u$ be an endomorphism of a bounded complex, in positive degrees, of
projective $A$-modules of finite type. One has \
\begin{equation*}\tag{1.8.1}\label{III:eq:1-8-1}
  \det(1-u t,K)^{-1} = \sum_n \tr(u,\Gamma^n K) t^n \text{.}
\end{equation*}
\end{proposition_}

Let $D(u,K)$ denote the right-hand side of \eqref{III:eq:1-8-1}.




\begin{lemma_}\label{III:1-9}
Let $u$ be an endomorphism of a short exact sequence of bounded complexes in
degrees $\geqslant 0$ of projective $A$-modules of finite type
\[
  0 \to K' \to K \to K'' \to 0 \text{.}
\]
One has $D(u,K)=D(u,K')\cdot D(u,K'')$.
\end{lemma_}

If $0\to M'\to M\to M''\to 0$ is a split short exact sequence of
$A$-modules, $\Gamma^n M$ has a natural filtration, with successive quotients
the $\Gamma^i M'\otimes \Gamma^j M''$ ($i+j=n$). Hence, $\Gamma^n K$ admits
a natural filtration, with successive quotients the complexes
$N(N^{-1} \Gamma^i K'\otimes N^{-1}\Gamma^j K'')$ ($i+j=n$), canonically
homotopic to the simple complexes $s(\Gamma^i K'\otimes \Gamma^i K'')$ deduced
from the double complexes $\Gamma^i K'\otimes \Gamma^i K''$, and
\[
  \tr(u,\Gamma^n K) = \sum_{i+j=n} \tr(u,\Gamma^i K')\cdot \tr(u,\Gamma^j K'') \text{,}
\]
whence the lemma. From \ref{III:1-9}, one deduces by induction.




\begin{lemma_}\label{III:1-10}
Let $F$ be a finite filtration of $K$ stable under $u$, such that the
$\gr_F^p(K^q)$ are projective of finite type. Then,
\[
  D(u,K) = \prod_p D\left(u,\gr_F^p(K)\right) \text{.}
\]
\end{lemma_}




\begin{lemma_}\label{III:1-11}
Let $u$ be an endomorphism of a projective $A$-module of finite type $M$, and
$M[-q]$ the complex reduced to $M$ in degree $q$. One has
\[
  D\left(u,M[-q]\right) = \left(\sum \tr(u,\Gamma^n M) t^n\right)^{(-1)^q} \qquad \text{($q\geqslant 0$.)}
\]
\end{lemma_}




By \ref{III:1-5}, this is true for $q=0$. Proceeding by induction, it
remains to show that
\[
  D\left(u,M[-q]\right) \cdot D\left(u,M[-(q+1)]\right) = 1 \text{.}
\]
This formula follows from \ref{III:1-9}, applied to the complex $K$ reduced to
$M$ in degrees $q$ and $q+1$ (with $d=\operatorname{id}$), and to its
stupid filtration:
\[
  (\cdots 0 \to M \cdots) \to (\cdots M \to M) \to (\cdots M \to 0 \to 0 \cdots) \text{.}
\]
Indeed, $K$ is homotopic to $0$. In the derived category, one has therefore
$\Gamma^n(K)=\Gamma^n(0)$; $0$ for $n>0$ and $A$ for $n=0$, and $D(u,K)=1$.




\subsection{}\label{III:1-12}

Let us prove \ref{III:1-8}. By \ref{III:1-10} applied to the
``stupid'' filtration of $K$ by the subcomplexes
\[
  0 \to \cdots \to 0 \to K^q \to K^{q+1} \to \cdots \text{,}
\]
one has
\[
  D(u,K) = \prod_q D\left(u,K^q[-q]\right) =_{\text{(\ref{III:1-10})}} \prod_q \left(\sum \tr\left(u,\Gamma^n(K^q)\right) t^n \right)^{(-1)^q} \text{.}
\]
It remains to prove that for $u$ an endomorphism of a projective module of finite
type $M$, one has
\begin{equation*}\tag{1.12.1}\label{III:eq:1-12-1}
  \det\left(1-u t,M\right)^{-1} = \sum \tr(u,\Gamma^n M) t^n \text{.}
\end{equation*}
If $u=0$, both sides equal $1$. Adding to $M$ a module $N$ on which
$u=0$, and applying \ref{III:1-9}, we reduce to assuming that $M$ is free.
Take a basis of it. It is then a matter of proving algebraic identities
between the coordinates of $u$. The principle of extension of algebraic
identities reduces us to assuming that $A$ is an algebraically closed
field. The module $M$ then admits a filtration stable under $u$ with successive
quotients of rank $1$, and \ref{III:1-10} reduces us to the case where $M$ is of
rank $1$. Formula \eqref{III:eq:1-12-1} then reduces to the identity
\[
  \frac{1}{1-a t} = \sum a^n t^n \text{.}
\]




\begin{corollary_}\label{III:1-13}
Let $u$ and $u'$ be endomorphisms of bounded complexes $K$ and $K'$ of
projective $A$-modules of finite type. If, in the derived category, $K$,
equipped with $u$, is isomorphic to $K'$, equipped with $u'$, then
\[
  \det(1-u t,K) = \det(1-u' t,K') \text{.}
\]
\end{corollary_}

Applying \ref{III:1-7-2}, we reduce to assuming that $K$ and $K'$ are in
positive degrees. One then observes that the right-hand side of \ref{III:1-8}
has the required invariance property.




\subsection{}\label{III:1-14}

This corollary makes it possible to define $\det(1-u t,K)$ for $u$ an endomorphism of
$K\in \ob\eD_\text{parf}(A)$. This construction is ad hoc; in fact, for $v$ an
automorphism of $L\in\ob\eD_\text{parf}(\Lambda)$, one can define
$\det(v)\in\Lambda^\times$, and $\det(1-u t,K)$ is obtained for $v=1-u t$,
$\Lambda=A\llbracket t\rrbracket$ and $L=K\otimes_A\Lambda$.










\section{The Theorem}\label{III:2}




\subsection{}\label{III:2-1}

Let us resume the notations of (Rapport, \S\ref{II:1}). For $A$ a
noetherian commutative torsion ring, $X_0$ a separated scheme of finite type over
$\dF_q$ and $\sK_0\in\ob\eD_\text{ctf}(X_0;A)$, we put
\[
  L(X_0,\sK_0) = \prod_{x\in |X_0|} \det\left(1-F_x^* t^{\deg(x)}, \sK\right)^{-1}
\]
(\ref{III:1-7}, \ref{III:1-13} et Rapport, \ref{II:1-5}, \ref{II:1-6}).




\begin{theorem_}\label{III:2-2}
In each of the following two cases
\begin{enumerate}[\indent a)]
  \item $A$ is torsion prime to $p$;
  \item $A$ is reduced and of characteristic $p$
\end{enumerate}
one has
\begin{equation*}\tag{2.2.1}\label{III:eq:2-2-1}
  L(X_0,\sK_0) = \det\left(1-F^* t^f, \eR\Gamma_c(X,\sK)\right)^{-1} \text{.}
\end{equation*}
\end{theorem_}




\subsection{Remark}\label{III:2-3}

In case a), the method of Rapport, \S\ref{II:3} makes it possible to deduce from
the trace formula (Rapport \ref{II:4-10}) that the logarithmic derivatives
of the two sides of \eqref{III:eq:2-2-1} are equal. The ring $A$ being of
torsion, this does not suffice to prove \ref{III:2-2}.




\subsection{Remark}\label{III:2-4}

If $A$ is a field, one has
\[
  L(X_0,\sK_0) = \prod_i L\left(X_0,\underline\h^i(\sK_0)\right)^{(-1)^i}
\]
and the spectral sequence
$E_2^{p q} = \h_c^p(X,\underline\h^q(\sK))\Rightarrow \h_c^{p+q}(X,\sK)$ shows
that
\begin{align*}
  \det\left(1-F^* t, \eR\Gamma_c(X,\sK)\right)
    &= \prod_n \det\left(1-F^* t^f,\h_c^n(X,\sK)\right)^{(-1)^n} \\
    &= \prod_{p,q} \det\left(1-F^* t,\h_c^p(X,\underline\h^q(\sK))\right)^{(-1)^{p+1}} \text{.}
\end{align*}
This reduces \ref{III:2-2} (for $A$ a field) to the following special case.


\subsubsection{}\label{III:2-4-1}

If $\sG_0$ is a constructible sheaf of $A$-vector spaces, one has
\[
  L(X_0,\sG_0) = \prod_n \det\left(1-F^* t^f,\h_c^n(X,\sG)\right)^{(-1)^{n+1}} \text{.}
\]




\subsection{Remark}\label{III:2-5}

It is easy to reduce case b) to that where $A$ is a field (or even
a finite field) of characteristic $p$, cf. \ref{III:4-2}.




\subsection{Remark}\label{III:2-6}

The special case of b) where $A=\dZ/p$, and where $X$ is reduced to the constant
sheaf $\dZ/p$ in degree $0$;
\[
  Z(X_0,t) = \prod_{x\in |X_0|} \left(1-t^{\deg(x)}\right)^{-1}
      \equiv \prod_n \det\left(1-F^* t^f,\h_c^n(X,\dZ/p)\right)^{(-1)^{n+1}} \pmod p
\]
had been proved by N. Katz \cite[XXII 3.1]{sga7}. His method, quite
different from the one followed here, starts from Dwork's $p$-adic theory of
$Z(X_0,t)$.




\subsection{Remark}\label{III:2-7}


Contrary to what I had imprudently claimed to friends, one cannot
in b) omit the hypothesis ``$A$ reduced.'' For a counter-
example, cf. \ref{III:4-5}.




\subsection{Plan of the proof of \ref{III:2-2}}\label{III:2-8}

A shift in degrees \ref{III:1-7-2} reduces us to assuming
that $\sK_0$ is a bounded complex, in positive degrees, of sheaves of
constructible and flat $A$-modules. One then has
$\eR\Gamma_c(X,\sK)\in \ob\eD_\text{parf}^{\geqslant 0}(A)$ (Rapport,
proof \ref{II:4-9}, b). We shall use \ref{III:1-8} to expand
both sides of \eqref{III:eq:2-2-1} as power series in $T=t^f$.
Equality of the coefficients of $T^n$ will follow from
\cite[XVII 5.5.21]{sga4} and from a trace formula on
$\operatorname{Sym}^n(X)$; in case a), the formula (Rapport,
\ref{II:4-10}), and in case b) a formula to be proved in the
next chapter. One should first reduce to the case where the
symmetric powers of $X_0$ exist (for example to the case where $X$ is
affine) using the multiplicativity properties in $X_0$ of the two
sides of \eqref{III:eq:2-2-1}: for $U_0$ an open subset of $X_0$, with closed
complement $Y_0$, one has $L(X_0,\sK_0)=L(U_0,\sK_0)\cdot L(Y_0,\sK_0)$ and
the analogous formula for the right-hand side follows from \ref{III:1-7-1}
by the method of (Rapport \S\ref{II:6}A).




\subsection{The case where $X_0$ is finite}\label{III:2-9}

The multiplicativity in $X_0$ of the two sides of \eqref{III:eq:2-2-1}
reduces this case to that where $X_0=\spec(\dF_{q^k})$. Returning to the
definitions, one sees that \eqref{III:eq:2-2-1} is then equivalent to
Rapport \ref{II:3-4}.




\subsection{The first member}\label{III:2-10}

By \ref{III:1-8}, one has
\[
  L(X_0,\sK_0) = \prod_{x\in |X_0|} \sum_n \tr\left(F_x^*,\Gamma^n\sK\right) t^{m\deg(x)} \text{.}
\]
The coefficient of $T^n$ is therefore the sum, extended to those of the
formal linear combinations $\sum_{x\in |X_0|} n_x\cdot x$
($n_x\geqslant 0$, almost all $n_x$ zero) of elements of $|X_0|$
such that $\sum n_x [k(x):\dF_q] = n$, of the products
\begin{equation*}\tag{2.10.1}\label{III:eq:2-10-1}
  \prod_x \tr\left(F_x^*,\Gamma^{n_x} \sK\right) \text{.}
\end{equation*}

For $x\in |X_0|$, let $(x)$ denote the $0$-cycle of the $[k(x):\dF_q]$ points of $X$
above $x$ (an orbit of $F$). The map
$\sum n_x\mapsto \sum n_x\cdot (x)$ identifies the formal linear
combinations as above, with $\sum n_x [k(x):\dF_q]=n$, with the effective
$0$-cycles of degree $n$ of $X$ fixed under $F$, i.e.\ with the fixed points of
$F$ in $\operatorname{Sym}^n(X)$. The coefficient of $T^n$ in $L(X_0,\sK_0)$
thus appears as a sum of products \eqref{III:eq:2-10-1}, indexed
by $\operatorname{Sym}^n(X)^F$. More precisely,




\begin{lemma_}\label{III:2-11}
 $L(X_0,\sK_0) = \sum_n \sum_{x\in\operatorname{Sym}^n(X)^F} \tr\left(F_x^*,\eL\Gamma_\textnormal{ext}^n \sK\right) T^n$.
\end{lemma_}

For $Y_0$ a larger and larger finite subscheme of $X_0$, one has
\[
  L(X_0,\sK_0) = \lim L(Y_0,\sK_0) \text{,}
\]
i.e., for every $m$, there exists a finite subscheme $Z_0\subset X_0$ such that
if $Y_0\supset Z_0$, $L(X_0,\sK_0)$ and $L(Y_0,\sK_0|Y_0)$ are congruent
$\mod T^m$.
It suffices that $Z_0$ contain the closed points of degree $\leqslant m$.
The same fact holding for the right-hand side, it suffices to prove
\ref{III:2-11} for $X_0$ finite. One then has
\begin{align*}
  L(X_0,\sK_0)
    &=_\text{\ref{III:2-9}} \det\left(1-F^* T,\Gamma(X,\sK)\right)^{-1} \\
    &=_\text{\ref{III:1-8}} \sum_n \tr\left(F^*,\Gamma^n \Gamma(X,\sK)\right) T^n \\
    &=_\text{\ref{III:1-3}} \sum_n \tr\left(F^*,\Gamma(\operatorname{Sym}^n(X),\Gamma_\text{ext}^n(\sK)\right) T^n \\
    &= \sum_n \sum_{x\in \operatorname{Sym}^n(X)^F} \tr\left(F_x^*,\Gamma_\text{ext}^n(\sK)\right) T^n
\end{align*}
(by the trace formula in the trivial case of a finite set).




\subsection{The second member}\label{III:2-12}

By \ref{III:1-8} and \cite[XVII 5.5.12]{sga4}, one has
\begin{align*}
  \det\left(1-F^* T,\eR\Gamma_c(X,\sK)\right)^{-1}
    &= \sum_n \tr\left(F^*,\eL\Gamma^n\eR\Gamma_c(X,\sK)\right) T^n \\
    &= \sum_n \tr\left(F^*,\eR\Gamma_c(\operatorname{Sym}^n(X),\Gamma_\text{ext}^n\sK)\right) T^n \text{.}
\end{align*}




\subsection{}\label{III:2-13}

The theorem is therefore equivalent to the trace formulas
\[
  \tr\left(F^*,\eR\Gamma_c(\operatorname{Sym}^n(X,\Gamma_\text{ext}^n \sK)\right) = \sum_{x\in\operatorname{Sym}^n(X)^F} \tr\left(F_x^*,\Gamma_\text{ext}^n\sK\right) \text{.}
\]
In case a), they follow from Rapport \ref{II:4-10}. The analogous
formula in case b) will be proved in \S\ref{III:4}.










\section{Artin-Schreier theory}\label{III:3}




\subsection{}\label{III:3-1}

In what follows, a coherent sheaf on a scheme $S$ will always be
regarded as a sheaf on $S_\text{\'et}$. Let $S$ be a scheme of
characteristic $p>0$. If $G$ is a locally constant sheaf of
$\dZ/p$-vector spaces of finite rank, $\sG=G\otimes_{\dZ/p}\sO$ is a
locally free coherent sheaf on $S$. We shall denote by $\Phi$ both the endomorphism
$f\mapsto f^p$ of $\sO$ and its tensor product with the identity of
$G$: $\phi:\sG\to\sG$. Artin-Schreier theory \cite[IX 3.5]{sga4}
asserts that the sequence
\[\xymatrix{
  0 \ar[r]
    & \dZ/p \ar[r]
    & \sO \ar[r]^-{\Phi-1}
    & \sO \ar[r]
    & 0
}\]
is exact. Tensoring with $G$, one finds an exact sequence
\[\xymatrix{
  0 \ar[r]
    & G \ar[r]
    & \sG \ar[r]^-{\Phi-1}
    & \sG \ar[r]
    & 0
}\]




\subsection{}\label{III:3-2}

Suppose that $S$ is an open subset in a noetherian scheme $\bar S$.
Let $j:S\hookrightarrow \bar S$ be the inclusion morphism, and $I$ a
sheaf of ideals defining the closed complement
$\bar S\setminus S$. Let us still denote by $\Phi$ the extension by functoriality
of $\Phi$ to $j_*\sG$.




\begin{lemma_}\label{III:3-3}
There exist coherent extensions $\bar\sG\subset j_*\sG$ of $\sG$ such
that
\begin{equation*}\tag{3.3.1}\label{III:eq:3-3-1}
  \Phi(\bar\sG)\subset I\cdot \bar\sG\text{.}
\end{equation*}
If $\bar\sG$ is such an extension, the sequence
\begin{equation*}\tag{3.3.2}\label{III:eq:3-3-2}
\xymatrix{
  0 \ar[r]
    & j_! G \ar[r]
    & \bar\sG \ar[r]^-{\Phi-1}
    & \bar\sG \ar[r]
    & 0
}
\end{equation*}
is exact.
\end{lemma_}
\begin{enumerate}[a)]
  \item \emph{Existence}: Let $\sG'\subset j_*\sG$ be a coherent extension
    of $\sG$ to $\bar S$. For every $n$, one has
    $\Phi(I^n \sG')\subset I^{p^n}\Phi(\sG')$. Let $n$ be large enough so
    that $I^{n(p-1)-1}\Phi(\sG')\subset \sG'$, and put $\bar\sG=I^n\sG'$.
    One has $\Phi(\bar\sG)\subset I^{p n}\Phi(\sG')
    =I\cdot I^n\cdot I^{n(p-1)-1}\Phi(\sG')\subset I\cdot I^n\sG'
    = I\bar\sG$.
  \item \emph{Kernel of $\Phi-1$}: Over $S$, the assertion follows from
    \ref{III:3-1}; it remains to check that a section $x$ of
    $\ker(\Phi-1)$, over $U$ \'etale over $\bar S$, vanishes in a neighbourhood of
    the inverse image of $\bar S\setminus S$. Since $x=\Phi(x)$, one has in
    fact by \eqref{III:eq:3-3-1} $x\in I\bar\sG$; moreover,
    $x\in I^n\bar\sG\Rightarrow x=\Phi(x)\in\Phi(I^n\bar\sG)\subset I^{n p}\bar\sG$, and $x$ lies in all the $I^n\bar\sG$, hence is zero in a neighbourhood of
    $\bar S\setminus S$.
  \item \emph{Surjectivity of $\Phi-1$}: One may assume $\bar S$ affine,
    write $\sG$ as a quotient of $\sO^n$; $\Phi$ then lifts to a
    $p$-linear map $\widetilde\Phi:\sO^n\to\sO^n$, and it suffices to
    check surjectivity of $\widetilde\Phi-1$:
    \[\xymatrix{
      \sO^n \ar[r]^-{\widetilde \Phi-1} \ar[d]
        & \sO^n \ar[d] \\
      \sG \ar[r]^-{\Phi-1}
        & \sG
    }\]
    The sheaf $\sO^n$ is the sheaf of local sections of the group
    scheme $\dG_a^n$, and $\widetilde\Phi-1$ is induced by a homomorphism
    of $\dG_a^n$ into $\dG_a^n$. This homomorphism is \'etale. To check
    that it is surjective, it suffices to see it after any base change
    $\spec(k)\to\bar S$ ($k$ algebraically closed) and over $k$ algebraically
    closed, every \'etale homomorphism between connected groups of the same dimension
    is surjective. Finally, an \'etale surjective morphism of
    $\bar S$-schemes induces an epimorphism between the corresponding
    sheaves on $\et{\bar S}$.
\end{enumerate}




\subsection{Frobenius and Frobenius}\label{III:3-4}

The absolute Frobenius morphism $F_\text{abs}:S\to S$ is the identity on
the topological space underlying $S$, and $f\mapsto f^p$ on the structural
sheaf. For $U/S$ \'etale, one has a canonical isomorphism
$F_\text{abs}^* U=U$: the diagram
\[\xymatrix{
  U \ar[r]^-{F_\text{abs}} \ar[d]
    & S \ar[d] \\
  S \ar[r]^-{F_\text{abs}}
    & S
}\]
is cartesian. The endomorphism of ringed site
$(\et{\bar S},\sO)\to(\et{\bar S},\sO)$ defined by $F_\text{abs}$ can therefore
be described as follows:
\begin{center}
  $F_\text{abs}^*$ is the identity on $\et S$; the homomorphism
  $F_\text{abs}^*\sO=\sO\to\sO$ is $\Phi$.
\end{center}
With this description, for every \'etale sheaf $\sG$ on $S$, the
Frobenius correspondence $F_\text{abs}^*\sG\to\sG$ is the identity.

Giving a $p$-linear morphism of quasi-coherent sheaves
$u:\sM\to\sN$ on $S$ amounts to giving a linear morphism
$u':F_\text{abs}^*\sM\to\sN$. In particular, with the notations of
\ref{III:3-5}, the homomorphism $\Phi:\sG\to\sG$ corresponds to
$\Phi':F_\text{abs}^*\sG\to\sG$.




\begin{lemma_}\label{III:3-5}
With the notations of \ref{III:3-4}, the diagram
\[\xymatrix{
  F_\textnormal{abs}^* G \ar[r] \ar[d]
    & G \ar[d] \\
  F_\textnormal{abs}^*\sG \ar[r]^-{\Phi'}
    & \sG
}\]
is commutative.
\end{lemma_}

It suffices to check it locally. One may therefore assume $G$ constant,
in which case the lemma is obvious.

The same theory holds for powers of $F_\text{abs}$.




\subsection{}\label{III:3-6}

Suppose now that $\bar S$ is a scheme over $\dF_q$ ($q=p^f$) and
in accordance with the conventions of Rapport \S\ref{II:1}, write
$S_0,\bar S_0,G_0,\sG_0,\bar\sG_0$ instead of $S,\dots$. Lemma
\ref{III:3-5}, for $F_\text{abs}^F:S_0\to S_0$, gives a commutative diagram
\[\xymatrix{
  F^* G_0 \ar[r] \ar[d]
    & G_0 \ar[d] \\
  F^*\sG_0 \ar[r]^-{(\Phi^f)'} \ar[r]
    & \sG_0
}\qquad\text{, then}\qquad
\xymatrix{
  F^*(j_! G_0) \ar[r] \ar[d]
    & j_!G_0 \ar[d] \\
  F^*\bar\sG_0 \ar[r]^-{(\Phi^f)'}
    & \sG_0
}\text{,}
\]
then, by scalar extension from $\dF_q$ to $\dF$ and passage to
cohomology
\[\xymatrix{
  F^*(j_! G) \ar[r]^-{F_*} \ar[d]
    & j_!G \ar[d] \\
  F^*\bar\sG \ar[r]
    & \bar\sG
}\qquad\text{and}\qquad
\xymatrix{
  \h^\bullet(\bar S,j_! G) \ar[r]^-{F^*} \ar[d]
    & \h^\bullet(\bar S,j_!G) \ar[d] \\
  \h^\bullet(\bar S,\bar\sG) \ar[r]
    & \h^\bullet(\bar S,\bar\sG)
}\text{,}
\]
where the second row is obtained by scalar extension from $\dF_q$ to
$\dF$ of the $\dF_q$-linear endomorphism $\Phi^f$ of
$\h^\bullet(\bar S_0,\bar G_0)$.

Artin-Schreier theory provides a long exact sequence
\[\xymatrix{
  \cdots \ar[r]
    & \h^i(\bar S,j_! G) \ar[r]
    & \h^i(\bar S,\bar\sG) \ar[r]^-{\Phi-1}
    & \h^i(\bar S,\bar\sG) \ar[r]
    & \cdots
}\]
where $\Phi$ is the $p$-linear extension to
$\h^\bullet(\bar S,\bar \sG)= \h^\bullet(\bar S_0,\bar\sG_0)\otimes_{\dF_q}\dF$
of the $p$-linear endomorphism $\Phi$ of $\h^\bullet(\bar S_0,\bar\sG_0)$.




\subsection{}\label{III:3-7}

If $\bar S_0$ is proper over $\dF_q$, the $\h^i(\bar S,\bar\sG)$ are
finite-dimensional and $\Phi-1$ is surjective (cf. \ref{III:3-6}b). One has therefore
\[
  \h_c^i(S,G) = \ker\left(\Phi-1 : \h^i(\bar S,\bar G) \to \h^i(\bar S,\bar\sG)\right) \text{,}
\]
and $F^*$ is induced by the $\dF$-linear extension to
$\h^i(\bar S,\bar\sG) = \h^i(\bar S_0,\bar\sG_0)\otimes_{\dF_q}\dF$ of
the $\dF_q$-linear endomorphism $\Phi^f$ of $\h^i(\bar S_0m\bar\sG_0)$.

Applying Lemma \ref{III:3-8} below, one deduces the identity
\begin{equation*}\tag{3.7.1}\label{III:eq:3-7-1}
  \tr\left(F^*,\h_c^i(S,G)\right) = \tr\left(\Phi^f,\h^i(\bar S_0,\bar\sG_0)\right) \text{.}
\end{equation*}




\begin{lemma_}\label{III:3-8}
Let $\Phi$ be a $p$-linear endomorphism of a vector space $V$ over
$\dF_q$: $\phi(\lambda x) = \lambda^p\phi(x)$. Put $F=\Phi^f$ ($F$ is
linear) and still denote by $\Phi$ and $F$ the respectively
$p$-linear and linear extensions of $\Phi$ and $F$ to $V\otimes_{\dF_q}\dF$. Then, $F$ stabilizes
$\ker(\Phi-1)\subset V\otimes_{\dF_q}\dF$ and
\[
  \tr\left(F,\ker(\Phi-1)\right) = \tr(F,V) = \tr\left(F,V\otimes_{\dF_q}\dF\right) \text{.}
\]
\end{lemma_}

On $V\otimes_{\dF_q}\dF$, one has $F\Phi = \Phi F$: the $p$-linear endomorphism
$F\Phi-\Phi F$ vanishes on $V$, hence everywhere. Write $V=V'+V''$, with $F$
invertible on $V'$ and nilpotent on $V''$. The theory of $p$-linear
endomorphisms ensures that
\[
  \ker(\Phi-1)\otimes_{\dF_q}\dF \iso V'\otimes_{\dF_q}\dF \text{.}
\]
Since $\tr(F,V'') = 0$, one has therefore
\begin{align*}
  \tr\left(F,\ker(\Phi-1)\right)
    &= \tr\left(F,\ker(\Phi-1)\otimes_{\dF_q}\dF\right) = \tr\left(F,V'\otimes_{\dF_q}\dF\right) \\
    &= \tr\left(F,V\otimes_{\dF_q}\dF\right) = \tr(F,V) \text{.}
\end{align*}










\section{Trace formula modulo \texorpdfstring{$p$}{p}}\label{III:4}




\begin{theorem_}\label{III:4-1}
Let $X_0$ be a separated scheme of finite type over $\dF_q$, $A$ a
commutative noetherian reduced ring of characteristic $p$, and
$\sK_0\in\ob\eD_\textnormal{ctf}(X_0,A)$. One has
\begin{equation*}\tag{4.1.1}\label{III:eq:4-1-1}
  \sum_{x\in X^F} \tr\left(F_x^*,\sK_x\right) = \tr\left(F^*,\eR\Gamma_c(X,\sK)\right) \text{.}
\end{equation*}
\end{theorem_}




\subsection{Reduction to the case where \texorpdfstring{$A$}{A} is a finite field}\label{III:4-2}

%NOTE: in this paragraph and others, (3.1.1) seems to be referring to (4.1.1)
A standard limit argument reduces us to the case where $A$ is
of finite type over $\dZ/p$. For each maximal ideal $\fm$ of $A$, the image of
each side of \eqref{III:eq:4-1-1} in $A/\fm$ is given by the same
formula, with $\sK_0$ replaced by
$\sK_0\lotimes_A A/\fm\in\ob\eD_\text{ctf}(X_0,A/\fm)$ (for the second side,
cf. Rapport \ref{II:4-12}). The $A/\fm$ being finite, and the intersection of the
maximal ideals of $A$ being reduced to $0$, one obtains the announced
reduction.




\subsection{Reduction to the case where \texorpdfstring{$A=\dZ/p$}{A=Z/p}}\label{III:4-3}

Suppose that $A$ is a finite field, and denote by $T_A'(\sK_0)$ and $T_A''(\sK_0)$
the two sides of \eqref{III:eq:4-1-1}. If, by restriction of scalars, one
regards $\sK_0$ as a complex of sheaves of $\dZ/p$-modules, one has
\[
  T_{\dZ/p}'(\sK_0) = \tr_{A/(\dZ/p)} T_A'(\sK_0)
\]
and similarly for $T''$. If \eqref{III:eq:4-1-1} holds for $A=\dZ/p$, one has therefore
\begin{equation*}\tag{4.3.1}\label{III:eq:4-3-1}
  \tr_{A/(\dZ/p)} T_A'(\sK_0) = \tr_{A/(\dZ/p)} T_A''(\sK_0) \text{.}
\end{equation*}
I claim that, for every $\lambda\in A$, one even has
\begin{equation*}\tag{4.3.2}\label{III:eq:4-3-2}
  \tr_{A/(\dZ/p)}\left(\lambda\cdot T_A'(\sK_0)\right) = \tr_{A/(\dZ/p)}\left(\lambda\cdot T_A''(\sK_0)\right) \text{,}
\end{equation*}
and hence $T_A'(\sK_0) = T_A''(\sK_0)$. This is clear for $\lambda=0$. For
$\lambda$ invertible, let $A(\lambda)$ denote the inverse image on $X$ of the
free rank-one $A$-module sheaf on $\spec(\dF_q)$ for which
$F^*=\lambda$. Formula \eqref{III:eq:4-3-2} is none other than
\eqref{III:eq:4-3-1} for $A(\lambda)\otimes_A \sK_0$.




% originally theorem 4.1bis (i.e. theorem 4.1 again)
\begin{theorem_}\label{III:4-1bis}
Let $\bar S_0$ be a smooth projective curve,
$j:S_0\hookrightarrow \bar S_0$ a dense open subset of $\bar S_0$, and $G_0$ a
locally constant sheaf of finite-rank $\dZ/p$-vector spaces on
$S_0$. One has
\begin{equation*}\tag{4.4.1}\label{III:eq:4-4-1}
  \sum_{x\in S^F} \tr\left(F_x^*,G_x\right) = \sum_i (-1)^i \tr\left(F^*,\h_c^i(S,G)\right) \text{.}
\end{equation*}
\end{theorem_}
\begin{proof}
Apply \ref{III:3-3} to $(S_0,\bar S_0,G_0)$. The coherent extension
$\bar\sG_0$ of $\sG_0$ whose existence \ref{III:3-3} guarantees is locally
free because it is torsion-free (since $\bar\sG_0\subset j_*\sG_0$) and
$\bar S_0$ is regular of dimension one. The $q$-linear endomorphism
$\Phi^f$ of $\bar\sG_0$ can be interpreted as a morphism
\[
  u : F^*\bar\sG_0 \to \bar\sG_0 \text{.}
\]
It factors through $I\bar G$, for $I$ defining the closed subset
$\bar S\setminus S$. Let $v:j_!G_0 \to \bar\sG_0$ be the unique extension of
the natural map $G_0 \to \sG_0 = G_0\otimes_{\dZ/p}\sO$. The diagram
\begin{equation*}\tag{4.4.2}\label{III:eq:4-4-2}
\xymatrix{
  F^* j_! G_0 \ar[r]^-{F^*} \ar[d]^-{F^* v}
    & j_! G_0 \ar[d]^-v \\
  F^*\bar \sG_0 \ar[r]^-u
    & \bar\sG_0
}
\end{equation*}
Finally, let $u F^*$ denote the composite map
\[\xymatrix{
  \h^\bullet(\bar S,\bar\sG) \ar[r]^-{F^*}
    & \h^\bullet(\bar S,F^*\bar\sG) \ar[r]^-u
    & \h^\bullet(\bar S,\bar\sG) \text{,}
}\]
one has by \eqref{III:eq:3-7-1}
\begin{equation*}\tag{4.4.3}\label{III:eq:4-4-3}
  \tr\left(F^*,\h_c^i(S,G)\right) = \tr\left(u F^*,\h^i(\bar S,\bar\sG)\right) \text{.}
\end{equation*}

%NOTE: trace down reference for Woodshole seminar
For each closed point $i_x:x\hookrightarrow \bar  S$ of $\bar S$, put
$\bar\sG_x = i_x^*\bar\sG$. The Lefschetz trace formula in coherent
cohomology (``Woodshole fixed point formula''), applied to $F$ and $u$, says
that
\begin{equation*}\tag{4.4.4}\label{III:eq:4-4-4}
  \sum_i (-1)^i \tr\left(u^* F^*, \h^i(\bar S,\bar\sG)\right) = \sum_{x\in \bar S^F} \frac{\tr(u_x,\sG_x)}{\det(1-d F_x)} \text{.}
\end{equation*}
On the right-hand side
\begin{enumerate}[\indent a)]
  \item since $d F=0$, the denominators are $1$;
  \item for $x\in \bar S^F\setminus S^F$, the endomorphism $u_x$ of $\bar\sG_x$
    is zero, hence so is $\tr(u_x,\bar\sG_x)$;
  \item for $x\in S^F$, \eqref{III:eq:4-4-4} provides a contribution
    \[
      \tr(u_x,\bar\sG_x) = \tr(F_x^*,G_x) \text{.}
    \]
    Identity \eqref{III:eq:4-4-1} therefore follows from \eqref{III:eq:4-4-3}
    and \eqref{III:eq:4-4-4}.
\end{enumerate}
\end{proof}




\subsection{A counterexample}\label{III:4-5}

We show by an example that one cannot, in \ref{III:4-1}, omit
the hypothesis that $A$ be reduced.

Let $Y_0$ be the elliptic curve over $\dF_2$ for which the Frobenius
eigenvalues are $\pm \sqrt{-2}$. It has $3$ rational points, and is
supersingular. The involution $\sigma:x\mapsto -x$ has therefore only one fixed
point, the point $0$. The quotient of $Y_0$ by $\sigma$ is a projective line;
that of $Y_0\setminus 0$ by $\sigma$ is therefore an affine line $\dF_0^1$, and
$Y_0\setminus 0$ is an unramified double covering of $\dF_0^1$:
\[
  \pi:Y_0\setminus 0 \to \dA_0^1 \qquad \text{, $\pi=\pi\sigma$.}
\]

The sheaf $\pi_*\dZ/2$ is a free module of rank one over the algebra $A$ of
the two-element group $\{1,\sigma\}$ over $\dZ/2$. This algebra is
isomorphic to that of dual numbers over $\dZ/2$.

The sheaf of $A$-modules $\pi_*\dZ/2$ is our example:
\begin{enumerate}[\indent a)]
  \item One has
    $\h_c^\bullet(\dA^1,\pi_*\dZ/2) = \h_c^\bullet(Y\setminus 0,\dZ/2)$.
    Since $Y_0$ is supersingular, one has $\h^i(Y,\dZ/2)=0$ for $i>0$, and
    the long exact sequence
    \[
      \cdots \to \h_c^i(Y\setminus 0,\dZ/2) \to \h^i(Y,\dZ/2) \to \h^i(\{0\},\dZ/2) \to \cdots
    \]
    shows that $\h_c^\bullet(Y\setminus 0,\dZ/2)=0$. One has therefore
    $\h_c^\bullet(\dA^1,\pi_*\dZ/2)=0$ and
    \[
      \tr_A\left(F^*,\eR\Gamma_c(\dA^1,\pi_*\dZ/2)\right) = 0 \text{.}
    \]
  \item Since $Y_0\setminus 0$ has two rational points, at one of the
    rational points of $\dA_1^0$, Frobenius substitution is the identity, and
    at the other, it is $\sigma$. One has therefore
    \[
      \sum_{x\in {\dA^1}^F} \tr\left(F^*,\pi_*\dZ/2\right) = 1+\sigma \text{,}
    \]
    and $1+\sigma\ne 0$.
\end{enumerate}

The equation of the covering $\pi$ is $y^2-y=x^3$.
